\input common
\def\begthm#1#2{
\vbox\bgroup
{\headerfont {Theorem #1:} \quad #2}
\smallskip\hrule\smallskip
}

\def\endthm{
\egroup
\bigskip\bigskip
}

\begthm{3.1}{Differentiable Implies Continuous}
If $f$ is differentiable at $a$, then $f$ is continuous at $a$.
\endthm

\begthm{3.1a}{Not Continuous implies Not Differentiable}
If $f$ is not continuous at $a$, then $f$ is not differentiable
at $a$.
\endthm

\begthm{3.2}{Constant Rule}
If $c$ is a real number, then ${d \over dx}(c) = 0$.
\endthm

\begthm{3.3}{Power Rule}
If $n$ is a nonnegative integer, then ${d \over dx}(x^n) = nx^{n - 1}$
\endthm

\begthm{3.4}{Constant Multiple Rule}
If $f$ is differentiable at $x$ and $c$ is a constant, then
$$\
{d \over dx}(cf(x)) = c f'(x).
$$
\endthm

\begthm{3.5}{Sum Rule}
If $f$ and $g$ are differentiable at $x$, then

$$\
{d \over dx}(f(x) + g(x)) = f'(x) + g'(x).
$$
\endthm

\begthm{3.6}{The Derivative of $e^x$}
The function $f(x) = e^x$ is differentiable for all real numbers $x$,
and
$$\
{d \over dx}(e^x) = e^x
$$
\endthm

\begthm{3.7}{Product Rule}
If $f$ and $g$ are differentiable at $x$, then

$$\
{d \over dx}(f(x)g(x)) = f'(x)g(x) + f(x)g'(x).
$$
\endthm

\begthm{3.8}{Quotient Rule}
If $f$ and $g$ are differntiable at $x$ and $g(x) \ne 0$,
then the derivative of $f / g$ at $x$ exists and
$$\
{d \over dx}\bigg({f(x) \over g(x)}\bigg) = {
g(x)f'(x) - f(x)g'(x)
\over
(g(x))^2
}
$$
\endthm

\begthm{3.9}{Power Rule (general form)}
If $n$ is any real number, then
$$\
{d \over dx}(x^n) = nx^{n - 1}
$$
\endthm

\begthm{2.1}{Relationship between One-Sided and Two-Sided Limits}
Assume $f$ is defined for all $x$ near $a$ except possibly
at $a$. Then

$$\
\lim_{x \to a} f(x) = L
$$

if and only if

$$\
\lim_{x \to a^{-}} = L
$$

and

$$\
\lim_{x \to a^{+}} = L
$$
\endthm

\begthm{2.2}{Limits of a Linear Function}
Let $a$, $b$, and $m$ be real numbers. For linear functions
$f(x) = mx + b$,

$$\
\lim_{x \to a} {f(x)} = f(a) = ma + b.
$$
\endthm

\begthm{2.3a}{Limit Law: Sum}
$$\
\lim_{x \to a}{
(f(x) + g(x))
} =
\lim_{x \to a}{f(x)} +
\lim_{x \to a}{g(x)} 
$$
\endthm

\begthm{2.3b}{Limit Law: Difference}
$$\
\lim_{x \to a}{
(f(x) - g(x))
} =
\lim_{x \to a}{f(x)} -
\lim_{x \to a}{g(x)} 
$$
\endthm

\begthm{2.3c}{Limit Law: Constant multiple}
$$\
\lim_{x \to a}{
cf(x)
} =
c\lim_{x \to a}{f(x)}
$$
\endthm

\begthm{2.3d}{Limit Law: Product}
$$\
\lim_{x \to a}{
(f(x)g(x))
} =
\bigg(
\lim_{x \to a}{f(x)}
\bigg)
\bigg(
\lim_{x \to a}{g(x)} 
\bigg)
$$
\endthm

\begthm{2.3e}{Limit Law: Quotient}
$$\
\lim_{x \to a}{
\bigg(
{
f(x)
\over
g(x)
}
\bigg)
} =
{
{\lim_{x \to a}{f(x)}}
\over
{\lim_{x \to a}{g(x)}}
}
$$

provided

$$\
\lim_{x \to a}{g(x)} \ne 0
$$
\endthm

\begthm{2.3f}{Limit Law: Power}
$$\
\lim_{x \to a}{
(f(x))^n
}
=
\big(
\lim_{x \to a}{
(f(x))
}
\big)^n
$$
\endthm

\begthm{2.3g}{Limit Law: Root}
$$\
\lim_{x \to a}{
(f(x))^{1/n}
}
=
\big(
\lim_{x \to a}{
(f(x))
}
\big)^{1/n}
$$

provided $f(x) > 0$ for $x$ near $a$, if $n$ is even.
\endthm

\begthm{2.4}{Limits of Polynomial and Rational Functions}
Assume $p$ and $q$ are polynomials and $a$ is a constant.

\par a. Polynomial functions
$$\
\lim_{x \to a} p(x) = p(a)
$$


\par b. Rational functions

$$\
\lim_{x \to a} {p(x) \over q(x)} = {p(a) \over q(a)}
$$

provided $q(a) \ne 0$.
\endthm

\begthm{2.3h}{Limit laws for One-Sided Limits}
Laws 2.3a-2.3f hold with $\lim\limits_{x \to a}$ replaced with
$\lim\limits_{x \to a^+}$ or $\lim\limits_{x \to a^-}$. Law 2.3g is
modified as follows.

\par a.

$$\
\lim_{x \to a^+}(f(x))^{1/n} =
\bigg(
\lim_{x \to a^+} f(x)
\bigg)^{1/n}
$$

provided $f(x) \ge 0$, for $x$ near $a$ with $x > a$, if $n$ is even.

\par b.

$$\
\lim_{x \to a^-}(f(x))^{1/n} =
\bigg(
\lim_{x \to a^-} f(x)
\bigg)^{1/n}
$$

provided $f(x) \ge 0$, for $x$ near $a$ with $x < a$, if $n$ is even.

\endthm

\begthm{2.5}{The Squeeze Theorem}
Assume the functions $f$, $g$, and $h$ satisfy $f(x) \le g(x) \le h(x)$
for all values of $x$ near $a$, except possibly at $a$.
If $\lim\limits_{x \to a} f(x) = \lim\limits_{x \to a} h(x) = L$,
then $\lim\limits_{x \to a} g(x) = L$.
\endthm

\begthm{2.6}{Limits at Infinity of Powers and Polynomials}
Let $n$ be a positive integer and let $p$ be the polynomial

$p(x) = a_n x^n + a_{n - 1}x^{n - 1} + \cdots + a_2 x^2 +
a_1x + a_0$, where $a_n \ne 0$.
\par 1.
$\lim\limits_{x \to \pm\infty} x^n = \infty$ when $n$ is even.

\par 2.
$\lim\limits_{x \to \infty} x^n = \infty$
and
$\lim\limits_{x \to -\infty} x^n = -\infty$
when $n$ is odd.

\par 3.
$\lim\limits_{x \to \pm\infty} {1 \over x^n} =
\lim\limits_{x \to \pm\infty} {x^{-n}} = 0
$.

\par 4.
$\lim\limits_{x \to \pm\infty} p(x)=
\lim\limits_{x \to \pm\infty} a_nx^n= \pm\infty$,
depending on the degree of the polynomial and the sign
of the leading coefficient $a_n$.

\endthm

\begthm{2.7}{End Behavior and Asymptotes of Rational Functions}
Suppose $f(x) = {p(x) \over q(x)}$ is a rational function,
where
$$\
p(x) = a_mx^m + a_{m-1}x^{m-1} + \cdots + a_2x^2 + a_0
$$

and

$$\
q(x) = b_nx^n + b_{n-1}x^{n-1} + \cdots + b_2x^2 + b_0
$$

with $a_m \ne 0$ and $b_n \ne 0$.

\par {\bf a. Degree of numerator less than degree of denominator.}
If $m < n$, then $\lim\limits_{x \to \pm\infty} f(x) = 0$,
and $y = 0$ is a horizontal asymptote of $f$.

\par {\bf b. Degree of numerator equals degree of denominator.}
If $m = n$, then $\lim\limits_{x \to \pm\infty} f(x) = a_m/b_n$,
and $y = a_m/b_n$ is a horizontal asymptote of $f$.

\par {\bf c. Degree of numerator greater than degree of denominator.}
If $m > n$, then $\lim\limits_{x \to \pm\infty} f(x) = \infty$ or $-\infty$,
and $f$ has no horizontal asymptote.

\par {\bf d. Slant asymptote.} 
If $m = n + 1$, then $\lim\limits_{x \to \pm\infty} f(x) = \infty$
or $-\infty$, and $f$ has no horizontal asymptote, but $f$
has a slant asymptote.

\par {\bf e. Vertifcal asymptotes.} Assuming $f$ is in
reduced form ($p$ and $q$ share no common factors), vertical
asymptotes occur at the zeros of $q$.

\endthm

\bye
